自旋 \(1/2\) 粒子的哈密顿量为 \(H(t)=-(\hbar\omega_0/2)\boldsymbol\sigma\cdot\mathbf n(t)\)，其中 \(\omega_0>0\)，\(\mathbf n(t)=(\sin\theta\cos\Omega t,\sin\theta\sin\Omega t,\cos\theta)\)，\(0<\theta<\pi\)、\(\Omega>0\)。初态为 \(t=0\) 时沿 \(\mathbf n(0)\) 的正自旋本征态；忽略外界退相干。给出一个完整绕转周期 \(T=2\pi/\Omega\) 的演化。
(1) 构造绕 \(z\) 轴共转的幺正变换，求时间无关的有效哈密顿量。
(2) 求旋转框架的能级间隔与初态相对于有效场轴线的夹角。
(3) 求一周期后自旋相对于瞬时原方向翻转的精确概率，并给出绝热极限判据。
(4) 在绝热极限写出两本征态积累的几何相位，区别动力相位及实验可观测量。
